\documentclass[12pt, a4paper]{article}

% ── Packages ────────────────────────────────────────────────────────────────
\usepackage[margin=1in]{geometry}
\usepackage{amsmath, amssymb, amsthm}
\usepackage{booktabs}
\usepackage{array}
\usepackage{enumitem}
\usepackage{hyperref}
\usepackage{titlesec}
\usepackage{fancyhdr}
\usepackage{xcolor}
\usepackage{mdframed}
\usepackage{tcolorbox}
\usepackage{microtype}
\usepackage{parskip}

% ── Colours ──────────────────────────────────────────────────────────────────
\definecolor{accentblue}{RGB}{0, 70, 140}
\definecolor{lightgray}{RGB}{245, 245, 245}
\definecolor{warnorange}{RGB}{200, 80, 0}

% ── Section formatting ───────────────────────────────────────────────────────
\titleformat{\section}{\large\bfseries\color{accentblue}}{}{0em}
  {\thesection\quad}[\titlerule]
\titleformat{\subsection}{\normalsize\bfseries\color{accentblue}}{}{0em}
  {\thesubsection\quad}

% ── Header / Footer ──────────────────────────────────────────────────────────
\pagestyle{fancy}
\fancyhf{}
\rhead{\small CSE400 -- Probability \& Computing}
\lhead{\small Lecture 23 Scribe}
\cfoot{\thepage}
\renewcommand{\headrulewidth}{0.4pt}

% ── Theorem-like environments ────────────────────────────────────────────────
\tcbuselibrary{theorems, skins, breakable}

\newtcbtheorem[number within=section]{defn}{Definition}%
  {colback=blue!5, colframe=accentblue, fonttitle=\bfseries,
   breakable}{def}

\newtcbtheorem[number within=section]{thm}{Theorem}%
  {colback=green!5, colframe=green!50!black, fonttitle=\bfseries,
   breakable}{thm}

\newtcbtheorem[number within=section]{exmp}{Example}%
  {colback=orange!5, colframe=warnorange, fonttitle=\bfseries,
   breakable}{ex}

% ── Macros ────────────────────────────────────────────────────────────────────
\newcommand{\EX}[1]{E\!\left[#1\right]}
\newcommand{\MGF}{M_X(t)}

% ─────────────────────────────────────────────────────────────────────────────
\begin{document}

% ── Title block ───────────────────────────────────────────────────────────────
\begin{center}
  {\LARGE\bfseries\color{accentblue}
    Predicting the Impossible:\\[4pt]
    Tail Bounds and System Reliability}\\[10pt]
  {\large\bfseries Case Study: Amazon Black Friday Sale 2026}\\[8pt]
  \rule{0.6\linewidth}{0.8pt}\\[6pt]
  \begin{tabular}{rl}
    \textbf{Course:}      & CSE400 -- Probability \& Computing \\
    \textbf{Instructor:}  & Dr.\ Dhaval Patel                  \\
    \textbf{Institution:} & Ahmedabad University               \\
    \textbf{Lecture:}     & 23                                 \\
    \textbf{Date:}        & April 02, 2026                     \\
  \end{tabular}
\end{center}

\vspace{8pt}

% ── Table of Contents ─────────────────────────────────────────────────────────
\tableofcontents
\newpage

% ─────────────────────────────────────────────────────────────────────────────
\section{Introduction: Amazon Black Friday Sale and Uncertainty}
% ─────────────────────────────────────────────────────────────────────────────

The lecture introduces a real-world computing reliability problem using the Amazon Black Friday
sale. The expected timeline is:

\begin{itemize}[leftmargin=2em]
  \item \textbf{Early Deals:} Nov 20--21, 2026
  \item \textbf{Main Sale:} Nov 25--27, 2026
  \item \textbf{Cyber Monday:} Nov 30, 2026
\end{itemize}

During such events, massive traffic spikes occur. The goal of the system architect is to answer:

\begin{tcolorbox}[colback=lightgray, colframe=accentblue, title=Core Question]
  \textit{What is the probability that incoming requests exceed the system capacity?}
  \[
    P(X \geq k)
  \]
\end{tcolorbox}

This question leads to the study of \textbf{tail probabilities}.

\begin{defn}{System Failure Event}{sfe}
  A system failure occurs when incoming requests exceed the system capacity:
  \[
    X \geq k
  \]
  where $X$ is the number of incoming requests and $k$ is the system threshold
  (maximum sustainable load). This is called a \textbf{tail probability} because
  it corresponds to the right tail of the distribution.
\end{defn}

% ─────────────────────────────────────────────────────────────────────────────
\section{Modeling Datacenter Traffic with Probability}
% ─────────────────────────────────────────────────────────────────────────────

A datacenter receives jobs across multiple servers $S_1, S_2, S_3, S_4$. Job arrivals
occur randomly, and some servers may receive $\geq k$ jobs. This traffic behaviour is
modelled using a \textbf{Poisson process}.

\begin{defn}{Poisson Arrival Model}{pam}
  A Poisson random variable models the number of events occurring in a fixed interval:
  \[
    P(X = k) = \frac{e^{-\lambda}\,\lambda^{k}}{k!}
  \]
  where $\lambda$ is the average arrival rate and $k$ is the number of arrivals.
\end{defn}

\begin{exmp}{Traffic Burst}{tb}
  \textbf{Condition:} System receives $\geq k$ jobs within one hour.\\
  \textbf{Probability of interest:} $P(X \geq k)$. This is the \emph{tail probability}.
\end{exmp}

% ─────────────────────────────────────────────────────────────────────────────
\section{Understanding Tail Events}
% ─────────────────────────────────────────────────────────────────────────────

A simulated server traffic graph uses the following parameters:
\begin{itemize}[leftmargin=2em]
  \item Poisson arrival rate: $\lambda = 50$
  \item Crash threshold: $k = 60$
\end{itemize}

\begin{tcolorbox}[colback=red!5, colframe=red!60!black, title=Key Observation]
  \textbf{``The mean is safe. The tail is catastrophic.''}
\end{tcolorbox}

Even though the mean is within safe limits, rare spikes cause crashes.

\medskip
\textbf{Observed Metrics (simulation duration: 1000 minutes):}
\begin{itemize}[leftmargin=2em]
  \item Crashes observed: 87
  \item Crash rate: $\dfrac{87}{1000} = 8.7\%$
\end{itemize}

Before launching the system, engineers must compute $P(X \geq 60)$.

% ─────────────────────────────────────────────────────────────────────────────
\section{The Hierarchy of Probability Bounds}
% ─────────────────────────────────────────────────────────────────────────────

\begin{center}
\begin{tabular}{cll}
  \toprule
  \textbf{Level} & \textbf{Bound}    & \textbf{Required Information} \\
  \midrule
  1              & Markov            & Mean $(\mu)$                  \\
  2              & Chebyshev         & Mean $+$ Variance $(\mu,\sigma^2)$ \\
  3              & Chernoff          & Moment Generating Function (MGF)   \\
  \bottomrule
\end{tabular}
\end{center}

\begin{tcolorbox}[colback=lightgray, colframe=accentblue]
  \textit{The more information we know about a distribution, the tighter the bound becomes.}
\end{tcolorbox}

% ─────────────────────────────────────────────────────────────────────────────
\section{Markov's Inequality}
% ─────────────────────────────────────────────────────────────────────────────

\begin{thm}{Markov's Inequality}{markov}
  For any non-negative random variable $X$ and $a > 0$:
  \[
    P(X \geq a) \leq \frac{\EX{X}}{a}
  \]
\end{thm}

This inequality provides an upper bound on tail probability using \emph{only} the mean.

\begin{exmp}{System Load}{sl}
  \textbf{Given:} Mean load $\EX{X} = 50$, crash threshold $a = 60$.

  \textbf{Applying Markov's inequality:}
  \[
    P(X \geq 60) \leq \frac{50}{60} \approx 0.833
  \]
  The probability of failure is \textbf{at most 83.3\%}.
\end{exmp}

\medskip
\textbf{Comparison:}
\begin{center}
\begin{tabular}{ll}
  \toprule
  \textbf{Method}  & \textbf{Probability} \\
  \midrule
  Exact            & 0.004  \\
  Markov estimate  & 0.714  \\
  \bottomrule
\end{tabular}
\end{center}

Markov overestimates heavily because it uses only the mean.

% ─────────────────────────────────────────────────────────────────────────────
\section{Chebyshev's Inequality}
% ─────────────────────────────────────────────────────────────────────────────

\begin{thm}{Chebyshev's Inequality}{cheb}
  For any random variable with variance $\sigma^2$:
  \[
    P\!\left(|X - \mu| \geq a\right) \leq \frac{\sigma^2}{a^2}
  \]
  \textbf{Tail form:}
  \[
    P(X \geq \mu + t) \leq \frac{\sigma^2}{t^2}
  \]
\end{thm}

\begin{exmp}{Chebyshev Applied to Server Load}{chebex}
  \textbf{Parameters:} $\mu = 50$, $k = 60$, $t = k - \mu = 10$.

  \textbf{Applying Chebyshev:}
  \[
    P(X \geq 60) \leq \frac{\sigma^2}{100}
  \]
  \textbf{Chebyshev estimate:} $0.125$
\end{exmp}

\medskip
\textbf{Comparison:}
\begin{center}
\begin{tabular}{ll}
  \toprule
  \textbf{Method}    & \textbf{Probability} \\
  \midrule
  Exact              & 0.004  \\
  Chebyshev estimate & 0.125  \\
  \bottomrule
\end{tabular}
\end{center}

Chebyshev is better than Markov but still loose.

% ─────────────────────────────────────────────────────────────────────────────
\section{Chernoff Bound and Moment Generating Functions}
% ─────────────────────────────────────────────────────────────────────────────

The Chernoff bound uses the Moment Generating Function (MGF).

\begin{defn}{Moment Generating Function}{mgf}
  \[
    M_X(t) = \EX{e^{tX}}
  \]
\end{defn}

\begin{thm}{Chernoff Bound}{chernoff}
  \[
    P(X \geq a) \leq \min_{t > 0}\; e^{-ta}\,M_X(t)
  \]
  The Chernoff bound transforms the probability using exponential functions,
  producing \textbf{exponentially decaying} bounds.
\end{thm}

\medskip
\textbf{Comparison:}
\begin{center}
\begin{tabular}{ll}
  \toprule
  \textbf{Method} & \textbf{Probability} \\
  \midrule
  Exact           & 0.092 \\
  Chernoff        & 0.391 \\
  \bottomrule
\end{tabular}
\end{center}

Chernoff produces the \textbf{tightest bound} among the three methods.

% ─────────────────────────────────────────────────────────────────────────────
\section{Visualization of Tail Decay}
% ─────────────────────────────────────────────────────────────────────────────

Three types of decay are compared:

\begin{center}
\begin{tabular}{lll}
  \toprule
  \textbf{Bound}  & \textbf{Decay Type}  & \textbf{Rate}    \\
  \midrule
  Markov          & Linear decay         & $1/k$            \\
  Chebyshev       & Polynomial decay     & $1/k^2$          \\
  Chernoff        & Exponential decay    & $e^{-k}$         \\
  \bottomrule
\end{tabular}
\end{center}

\begin{tcolorbox}[colback=lightgray, colframe=accentblue, title=Key Insight]
  Exponential decay reaches near-zero probability \textbf{much faster} than linear or
  polynomial decay.
\end{tcolorbox}

% ─────────────────────────────────────────────────────────────────────────────
\section{Case Study: Black Friday SLA Breach}
% ─────────────────────────────────────────────────────────────────────────────

A full 24-hour sale simulation uses the following parameters:
\begin{itemize}[leftmargin=2em]
  \item Average traffic: $\lambda = 500$
  \item SLA crash threshold: $k = 600$
\end{itemize}

\textbf{Mission:} The architect must estimate $P(X \geq 600)$ \emph{before} launching the sale.

% ─────────────────────────────────────────────────────────────────────────────
\section{Cost of Loose Probability Bounds}
% ─────────────────────────────────────────────────────────────────────────────

\textbf{Probability estimates compared:}
\begin{center}
\begin{tabular}{ll}
  \toprule
  \textbf{Method} & \textbf{Probability} \\
  \midrule
  Exact           & 1.4\%   \\
  Chernoff        & 8.8\%   \\
  Chebyshev       & 20\%    \\
  Markov          & 90.9\%  \\
  \bottomrule
\end{tabular}
\end{center}

Markov suggests almost certain failure, which is \textbf{incorrect}. The waste multiplier
is defined as:
\[
  \text{Waste} = \frac{\text{Bound Estimate}}{\text{Exact Probability}}
\]

\textbf{Financial over-provisioning waste:}
\begin{center}
\begin{tabular}{ll}
  \toprule
  \textbf{Method} & \textbf{Waste Multiplier} \\
  \midrule
  Chernoff        & $10.7\times$        \\
  Chebyshev       & $6422\times$        \\
  Markov          & $107039\times$      \\
  \bottomrule
\end{tabular}
\end{center}

\begin{tcolorbox}[colback=red!5, colframe=red!60!black]
  \textbf{Poor probability bounds lead to massive financial waste.}
\end{tcolorbox}

% ─────────────────────────────────────────────────────────────────────────────
\section{Choosing the Correct Tail Bound}
% ─────────────────────────────────────────────────────────────────────────────

\begin{center}
\begin{tabular}{p{2.2cm} p{2.8cm} p{2cm} p{2.2cm} p{3.5cm}}
  \toprule
  \textbf{Method} & \textbf{Required Data} & \textbf{Tail Decay}
    & \textbf{Tightness} & \textbf{Use Case} \\
  \midrule
  Markov      & Mean                    & $1/k$     & Very loose & Rough estimate \\[4pt]
  Chebyshev   & Mean + Variance         & $1/k^2$   & Moderate   & Unknown distribution \\[4pt]
  Chernoff    & Full MGF                & $e^{-k}$  & Tight      & Known models (Poisson, Binomial) \\
  \bottomrule
\end{tabular}
\end{center}

% ─────────────────────────────────────────────────────────────────────────────
\section{Class Exercise: Bounding the Laplace Distribution}
% ─────────────────────────────────────────────────────────────────────────────

\textbf{Distribution:} $X \sim \mathrm{Laplace}(\mu, b)$ with $\sigma = b\sqrt{2}$.

\textbf{Properties:}
\begin{itemize}[leftmargin=2em]
  \item Sharp central peak
  \item Heavy tails
\end{itemize}

\textbf{Goal:} Apply the Chernoff bound with threshold $a = \mu + 3\sigma$.

\subsection*{Derivation Blueprint}

\begin{enumerate}[leftmargin=2em, label=\textbf{Step \arabic*.}]
  \item \textbf{Compute MGF:}
    \[
      M_X(t) = \EX{e^{tX}}
    \]
  \item \textbf{Apply Chernoff inequality:}
    \[
      P(X \geq a) \leq e^{-ta}\,M_X(t)
    \]
  \item \textbf{Optimise:} Take derivative with respect to $t$, set it to zero, and solve
        for the optimal $t^*$.
  \item \textbf{Substitute} the optimal $t^*$ back into the bound.
\end{enumerate}

% ─────────────────────────────────────────────────────────────────────────────
\section{Summary and Exam-Oriented Memory Guide}
% ─────────────────────────────────────────────────────────────────────────────

\subsection*{Key Concepts (Flashcard Style)}

\begin{itemize}[leftmargin=2em]
  \item \textbf{Tail Probability} $\longrightarrow$ $P(X \geq k)$
  \item \textbf{System Reliability} $\longrightarrow$ Avoid overload
  \item \textbf{Markov} $\longrightarrow$ Mean only
  \item \textbf{Chebyshev} $\longrightarrow$ Mean $+$ Variance
  \item \textbf{Chernoff} $\longrightarrow$ MGF
  \item \textbf{Decay speed:} Linear $<$ Polynomial $<$ Exponential
\end{itemize}

\subsection*{Important Formulas}

\begin{tcolorbox}[colback=lightgray, colframe=accentblue, title=Formula Sheet]
  \textbf{Markov:}
  \[
    P(X \geq a) \leq \frac{\EX{X}}{a}
  \]

  \textbf{Chebyshev:}
  \[
    P(|X - \mu| \geq k) \leq \frac{\sigma^2}{k^2}
  \]

  \textbf{Chernoff:}
  \[
    P(X \geq a) \leq \min_{t > 0}\; e^{-ta}\,M_X(t)
  \]

  \textbf{Poisson distribution:}
  \[
    P(X = k) = \frac{e^{-\lambda}\,\lambda^k}{k!}
  \]
\end{tcolorbox}

\subsection*{Applications in Systems}

\begin{itemize}[leftmargin=2em]
  \item Datacenter scaling
  \item Network congestion
  \item Cloud capacity planning
  \item SLA reliability prediction
  \item Cyber Monday traffic spikes
\end{itemize}

\subsection*{Formula Application Tricks}

\begin{itemize}[leftmargin=2em]
  \item \textbf{Markov} --- Use when only the mean is known.
  \item \textbf{Chebyshev} --- Use when variance is known but the distribution is unknown.
  \item \textbf{Chernoff} --- Use when the distribution is known (Poisson, Binomial).
\end{itemize}

\subsection*{Common Exam Mistakes}

\begin{enumerate}[leftmargin=2em]
  \item Using Markov when variance is available.
  \item Forgetting that Chernoff requires the MGF.
  \item Confusing $P(|X - \mu| \geq a)$ with $P(X \geq a)$.
  \item Forgetting to optimise parameter $t$ in the Chernoff bound.
  \item Misinterpreting bounds as exact probabilities.
\end{enumerate}

\end{document}